% Edited conversion of source p. 48.
\ReportHeading
  {Perturbed Motions of a Nearly Dynamically Spherical Gyrostat with a Moving Mass Subjected to Constant Body-Fixed Torques}
  {Dmytro Leshchenko, Tetiana Kozachenko}
  {Odesa State Academy of Civil Engineering and Architecture}
  {}

The motion of a rigid body about a fixed point is a classical problem of mechanics. Its importance increased substantially with the development of spacecraft and satellite dynamics, because vehicles rotating about their centres of mass are subjected to torques of different physical origins. These include torques generated by rotors and gyroscopes, crew motion, loosely attached components, and fluid motion in internal cavities.

The dynamics of a rigid body containing moving internal masses is therefore an important extension of the classical Euler problem. In the present model, a viscoelastic internal element is represented by a moving mass connected to the body through a strongly damped elastic linkage. Such a model can describe loosely attached components whose long-term motion affects the attitude dynamics of a space vehicle. The body may also contain a viscous fluid in a cavity and is subjected to constant torques fixed in the body frame.

The objective is to investigate the combined influence of the viscous fluid, the moving mass, and constant body-fixed torques on a nearly dynamically spherical gyrostat. This combination extends earlier studies in which the perturbations were considered separately.

The equations of motion are reduced to a standard nonlinear system. The averaging method is then applied to obtain a simpler system that describes the slow evolution of the motion over a long time interval. The averaged equations make it possible to establish qualitative and quantitative properties of the rotation and to determine how the angular motion evolves under the simultaneous action of the different perturbations.

The results are relevant to the analysis of artificial satellites, spacecraft attitude dynamics, crew-induced disturbances, moving-mass control systems, and re-entry vehicles. Further work will focus on constructing explicit asymptotic approximations to the Euler equations and on integrating the averaged system for perturbation torques of different physical origins.
